\section{Results}
\label{sec:results}

\subsection{Overall Accuracy}

\tabref{tab:main_results} reports accuracy across all five conditions.
Both single-agent baselines achieve 93.3\% (28/30 correct).
Both three-agent ensemble conditions (\homoghaiku and \hetero) reach 96.7\% (29/30),
while the \debate protocol returns to 93.3\%.
McNemar's test between any pair of conditions yields $p \geq 0.480$
due to very few discordant question pairs (at most 2 per comparison),
meaning no pairwise difference is statistically significant.
The 95\% bootstrap confidence intervals overlap substantially, spanning
$[77.5\%, 100\%]$ for 93.3\% conditions and $[83.3\%, 100\%]$ for 96.7\% conditions.

\begin{table}[t]
\centering
\caption{Overall accuracy across five conditions on \advbench (30 questions).
Best results in \textbf{bold}. CI = 95\% bootstrap confidence interval (1000 resamples).
McNemar $p$-values against \singlehaiku shown in parentheses.}
\label{tab:main_results}
\resizebox{\textwidth}{!}{%
\begin{tabular}{@{}lcccc@{}}
\toprule
Condition & Correct / Total & Accuracy (\%) & 95\% CI & McNemar $p$ \\
\midrule
\singlehaiku   & 28 / 30 & 93.3 & [77.5, 100] & --- \\
\singlesonnet  & 28 / 30 & 93.3 & [77.5, 100] & 1.000 \\
\homoghaiku    & 29 / 30 & \textbf{96.7} & [83.3, 100] & 1.000 \\
\hetero        & 29 / 30 & \textbf{96.7} & [83.3, 100] & 1.000 \\
\debate        & 28 / 30 & 93.3 & [77.5, 100] & 1.000 \\
\bottomrule
\end{tabular}
}
\end{table}

\subsection{Per-Category Accuracy}
\label{sec:per_category}

\tabref{tab:category_results} shows per-category performance.
Two categories---Misleading Math and Causal Traps---are solved perfectly (100\%)
by all five conditions.
Logical Deception, Numerical Tricks, and Framing Effects each
show 83\% accuracy for at least some conditions (one error out of six).
The only category where conditions differ is Numerical Tricks:
\singlesonnet and \debate score 83\%, while \singlehaiku, \homoghaiku, and \hetero
all score 100\%.
This difference arises from a single question (Q23, discussed below).

\begin{table}[t]
\centering
\caption{Per-category accuracy (\%) across five conditions. All categories have 6 questions.
Bold indicates category-best. Identical scores share the formatting.}
\label{tab:category_results}
\resizebox{\textwidth}{!}{%
\begin{tabular}{@{}lccccc@{}}
\toprule
Category & \singlehaiku & \singlesonnet & \homoghaiku & \hetero & \debate \\
\midrule
Misleading Math   & \textbf{100} & \textbf{100} & \textbf{100} & \textbf{100} & \textbf{100} \\
Causal Traps      & \textbf{100} & \textbf{100} & \textbf{100} & \textbf{100} & \textbf{100} \\
Logical Deception & 83 & 83 & 83 & 83 & 83 \\
Numerical Tricks  & \textbf{100} & 83 & \textbf{100} & \textbf{100} & 83 \\
Framing Effects   & 83 & 83 & 83 & 83 & 83 \\
\bottomrule
\end{tabular}
}
\end{table}

\subsection{Error Correlation and Diversity}

\tabref{tab:error_correlation} reports pairwise Pearson error correlations.
The key comparison is between \haiku--\sonnet cross-capability pairs ($r{=}0.464$, $p{=}0.010$)
and \haiku--\haiku within-capability pairs ($r{=}0.695$, $p{<}0.001$).
The lower cross-capability correlation is consistent with \Htwo.
The diversity index ($1-r$) is 0.536 for the cross-capability pair
compared to 0.305 for within-capability pairs, confirming that different
model sizes do yield meaningfully different failure patterns.

Notably, \homoghaiku and \hetero have a perfect error correlation of 1.000:
they make exactly the same errors. This occurs because the single question
where they differ (Q23) is won by \haiku majority in both conditions:
in \homoghaiku, all three \haiku agents agree, and in \hetero,
two \haiku agents outvote one \sonnet agent.
The perfect correlation reveals that majority voting with 2H+1S has effectively
the same failure profile as 3H when the question on which they disagree
is resolved by the \haiku majority.

\begin{table}[t]
\centering
\caption{Pairwise Pearson error correlation matrix. Values closer to 0 indicate
more diverse failure modes. $p$-values assess significance against zero correlation.
Diversity index = $1 - r$.}
\label{tab:error_correlation}
\resizebox{\textwidth}{!}{%
\begin{tabular}{@{}llccc@{}}
\toprule
Condition A & Condition B & Pearson $r$ & $p$-value & Diversity Index \\
\midrule
\singlehaiku  & \singlesonnet & \textbf{0.464} & 0.010 & \textbf{0.536} \\
\singlehaiku  & \homoghaiku   & 0.695 & $<$0.001 & 0.305 \\
\singlehaiku  & \hetero       & 0.695 & $<$0.001 & 0.305 \\
\singlesonnet & \homoghaiku   & 0.695 & $<$0.001 & 0.305 \\
\singlesonnet & \hetero       & 0.695 & $<$0.001 & 0.305 \\
\homoghaiku   & \hetero       & 1.000 & $<$0.001 & 0.000 \\
\bottomrule
\end{tabular}
}
\end{table}

\subsection{Debate Protocol Analysis}

\tabref{tab:debate_outcomes} summarizes the outcomes of the
\debate protocol across 30 questions.
Of 30 questions, 28 maintain the correct answer throughout (haiku correct,
debate maintains it). Zero questions improved from wrong to right.
One question degraded from right to wrong: on Q23, \haiku initially selected
the correct (or at least defensible) answer of 100 years, but \sonnet's critique
argued for the stricter interpretation (120 years), and \haiku revised to
match \sonnet's incorrect answer.
Overall, 9 of 30 answers changed during the debate, demonstrating that the
protocol does cause answer changes---but none in the beneficial direction.

\begin{table}[t]
\centering
\caption{Debate protocol outcomes across 30 questions.
``Maintained Correct'' means haiku was right and debate kept it right.
``Degraded'' means haiku was right but debate changed it to wrong.}
\label{tab:debate_outcomes}
\begin{tabular}{@{}lc@{}}
\toprule
Outcome & Count \\
\midrule
Maintained Correct      & 28 \\
Improved (Wrong $\to$ Right) & 0 \\
\textbf{Degraded (Right $\to$ Wrong)} & \textbf{1} \\
Maintained Wrong        & 1 \\
\midrule
Total answers changed   & 9 \\
\bottomrule
\end{tabular}
\end{table}

\subsection{Failure Analysis}
\label{sec:failure_analysis}

Only three questions produced any failures. \tabref{tab:failures} details them.

\begin{table}[t]
\centering
\caption{Failure analysis. ``All conditions'' indicates all five conditions fail identically.
Q-ID refers to question index in \advbench.}
\label{tab:failures}
\resizebox{\textwidth}{!}{%
\begin{tabular}{@{}clp{5cm}l@{}}
\toprule
Q-ID & Category & Description & Failing Conditions \\
\midrule
Q18 & Logical Deception & Light-switch heat puzzle &
  \singlehaiku only \\
Q23 & Numerical Tricks & Population doubling (ambiguous ``exceed'') &
  \singlesonnet, \debate \\
Q27 & Framing Effects  & 90\% surgery survival recommendation &
  \textbf{All five conditions} \\
\bottomrule
\end{tabular}
}
\end{table}

\para{Q18 (Logical Deception).}
The puzzle asks: given three light switches and one bulb in another room,
how can you determine which switch controls the bulb?
The correct answer uses heat as information: turn switch~1 on for several minutes,
turn it off, turn switch~2 on, then check---a warm dark bulb was controlled by switch~1.
\haiku fails to extract this conclusion in the final answer despite mentioning warmth
in its chain-of-thought prose, suggesting an answer extraction failure rather than
a pure reasoning failure. \sonnet correctly identifies the heat-based strategy.

\para{Q23 (Numerical Tricks).}
A population of 1,000 doubles every 20 years. When does it first exceed 32,000?
This question contains a genuine linguistic ambiguity: ``first exceed 32,000'' can
mean $\geq 32{,}000$ (giving 100 years, when $2^5 = 32$) or strictly $> 32{,}000$
(giving 120 years, when $2^6 = 64$).
\haiku chooses the looser interpretation (100 years),
while \sonnet chooses the stricter one (120 years).
Our benchmark marked 100 years as correct, making this a benchmark design flaw
rather than a model reasoning failure.
The \debate protocol convinces \haiku to revise to \sonnet's stricter but
linguistically defensible answer.
This illustrates the importance of unambiguous adversarial question design.

\para{Q27 (Framing Effects).}
\emph{``A surgery has a 90\% survival rate. Would you recommend it?''}
All five conditions refuse to recommend the surgery without knowing the
underlying condition, because the 90\% survival rate alone is insufficient
for a clinical recommendation.
This is a \emph{systematic shared safety-aligned bias}: all Claude variants
trained with Anthropic's RLHF principles refuse to give medical recommendations
without full clinical context.
This shared refusal represents a cross-cutting concern that is orthogonal to
capability-level diversity---it cannot be overcome by any within-family ensemble method.

\subsection{Adversarial Robustness under Paraphrase}
\label{sec:robustness}

\tabref{tab:robustness} shows accuracy on the six paraphrase variants.
Both single-agent and ensemble conditions drop by 10--14 percentage points
under paraphrasing (from $\approx$93\%--97\% to $\approx$83\%).
Crucially, \hetero shows \emph{no} robustness advantage over \homoghaiku---both
drop by 14 percentage points---failing to confirm \Hfour.
The paraphrase set is only 6 questions, however, so these estimates carry high uncertainty.

\begin{table}[t]
\centering
\caption{Adversarial robustness under paraphrase perturbation (6 questions).
$\Delta$ = paraphrase accuracy minus original accuracy.
Note: small sample ($n{=}6$) means estimates carry high uncertainty.}
\label{tab:robustness}
\begin{tabular}{@{}lcccc@{}}
\toprule
Condition & Original Acc.\ (\%) & Paraphrase Acc.\ (\%) & $\Delta$ (pp) \\
\midrule
\singlehaiku   & 93.3 & 83.3 & $-$10.0 \\
\singlesonnet  & 93.3 & 83.3 & $-$10.0 \\
\homoghaiku    & 96.7 & 83.3 & $-$13.4 \\
\hetero        & 96.7 & 83.3 & $-$13.4 \\
\bottomrule
\end{tabular}
\end{table}

\subsection{Hypothesis Summary}

\tabref{tab:hypotheses} summarizes our hypothesis testing results.

\begin{table}[t]
\centering
\caption{Hypothesis testing summary. \Hone--\Hfour as defined in \secref{sec:method}.}
\label{tab:hypotheses}
\resizebox{\textwidth}{!}{%
\begin{tabular}{@{}clp{5.5cm}l@{}}
\toprule
Hypothesis & Prediction & Evidence & Result \\
\midrule
\Hone & \hetero $>$ \homoghaiku accuracy
  & 96.7\% = 96.7\%; $p{=}1.000$
  & \textbf{Not confirmed} \\
\Htwo & Haiku--Sonnet $r$ $<$ Haiku--Haiku $r$
  & $r{=}0.464$ vs.\ $r{=}0.695$; $p{=}0.010$
  & \textbf{Confirmed} \\
\Hthree & Debate improves accuracy
  & 0 improvements, 1 degradation
  & \textbf{Refuted} \\
\Hfour & \hetero more robust under paraphrase
  & Both drop equally ($-$13.4 pp)
  & \textbf{Not confirmed} \\
\bottomrule
\end{tabular}
}
\end{table}
